\documentclass[10pt,a4paper]{amsart}
\usepackage{amsmath,amssymb,amsthm,graphicx}
% Retained preamble for the Stacks Project chapter morphisms.tex (Morphisms of Schemes).
% The chapter file pinned-original-source.tex carries only \input{preamble} on its first line, so no
% preamble text lives inside the chapter itself.  The Stacks Project uses one single shared preamble for
% every chapter; the snapshot stages it verbatim in sources/stacks-chow/preamble.tex.  The block below is
% copied mechanically, line for line, from that file's lines 14-18 (the Xy-pic setup) and lines 154-217
% (the native theorem declarations together with the mathematical macros).  Nothing was added, removed or
% redefined.  The wrapper class and the amsmath/amssymb/amsthm/graphicx base packages are supplied by the
% assembler, so the class/\externaldocument part of the shared preamble is not repeated here.
% For commutative diagrams we use Xy-pic
\usepackage[all]{xy}

% We use 2cell for 2-commutative diagrams.
\xyoption{2cell}
% Theorem environments.
%
\theoremstyle{plain}
\newtheorem{theorem}[subsection]{Theorem}
\newtheorem{proposition}[subsection]{Proposition}
\newtheorem{lemma}[subsection]{Lemma}

\theoremstyle{definition}
\newtheorem{definition}[subsection]{Definition}
\newtheorem{example}[subsection]{Example}
\newtheorem{exercise}[subsection]{Exercise}
\newtheorem{situation}[subsection]{Situation}

\theoremstyle{remark}
\newtheorem{remark}[subsection]{Remark}
\newtheorem{remarks}[subsection]{Remarks}

\numberwithin{equation}{subsection}

% Macros
%
\def\lim{\mathop{\mathrm{lim}}\nolimits}
\def\colim{\mathop{\mathrm{colim}}\nolimits}
\def\Spec{\mathop{\mathrm{Spec}}}
\def\Hom{\mathop{\mathrm{Hom}}\nolimits}
\def\Ext{\mathop{\mathrm{Ext}}\nolimits}
\def\SheafHom{\mathop{\mathcal{H}\!\mathit{om}}\nolimits}
\def\SheafExt{\mathop{\mathcal{E}\!\mathit{xt}}\nolimits}
\def\Sch{\mathit{Sch}}
\def\Mor{\mathop{\mathrm{Mor}}\nolimits}
\def\Ob{\mathop{\mathrm{Ob}}\nolimits}
\def\Sh{\mathop{\mathit{Sh}}\nolimits}
\def\NL{\mathop{N\!L}\nolimits}
\def\CH{\mathop{\mathrm{CH}}\nolimits}
\def\proetale{{pro\text{-}\acute{e}tale}}
\def\etale{{\acute{e}tale}}
\def\QCoh{\mathit{QCoh}}
\def\Ker{\mathop{\mathrm{Ker}}}
\def\Im{\mathop{\mathrm{Im}}}
\def\Coker{\mathop{\mathrm{Coker}}}
\def\Coim{\mathop{\mathrm{Coim}}}

% Boxtimes
%
\DeclareMathSymbol{\boxtimes}{\mathbin}{AMSa}{"02}

%
% Macros for moduli stacks/spaces
%
\def\QCohstack{\mathcal{QC}\!\mathit{oh}}
\def\Cohstack{\mathcal{C}\!\mathit{oh}}
\def\Spacesstack{\mathcal{S}\!\mathit{paces}}
\def\Quotfunctor{\mathrm{Quot}}
\def\Hilbfunctor{\mathrm{Hilb}}
\def\Curvesstack{\mathcal{C}\!\mathit{urves}}
\def\Polarizedstack{\mathcal{P}\!\mathit{olarized}}
\def\Complexesstack{\mathcal{C}\!\mathit{omplexes}}
% \Pic is the operator that assigns to X its picard group, usage \Pic(X)
% \Picardstack_{X/B} denotes the Picard stack of X over B
% \Picardfunctor_{X/B} denotes the Picard functor of X over B
\def\Pic{\mathop{\mathrm{Pic}}\nolimits}
\def\Picardstack{\mathcal{P}\!\mathit{ic}}
\def\Picardfunctor{\mathrm{Pic}}
\def\Deformationcategory{\mathcal{D}\!\mathit{ef}}

% --- appended 2026-10-12 (Chief): the project own definitions of the reference/slogan/history
% display environments, copied mechanically from the shared Stacks preamble lines 8-12
% (\usepackage{verbatim} plus the three \newenvironment lines, which the chapter preamble keeps
% verbatim from the project preamble). Nothing added, removed or redefined: these environments
% carry no mathematical content and the project itself typesets them as comments. Without these
% three lines a counted statement whose own text uses them fails with Environment slogan undefined.
\usepackage{verbatim}
\newenvironment{reference}{\comment}{\endcomment}
\newenvironment{slogan}{\comment}{\endcomment}
\newenvironment{history}{\comment}{\endcomment}
\usepackage{cleveref}
\crefname{originallemma}{original lemma}{original lemmas}
\Crefname{originallemma}{Original lemma}{Original lemmas}
\makeatletter
\newlabel{morphisms-lemma-finite-degree}{{02NX}{source}{}{sourcepointer-1}{}}
\newlabel{morphisms-lemma-finite-degree@cref}{{[originallemma][0][]02NX}{[1][0][]source}{}{}{}}
\makeatother
\crefname{originallemma}{original lemma}{original lemmas}
\Crefname{originallemma}{Original lemma}{Original lemmas}
\makeatletter
\newlabel{lemma-locally-finite-type-characterize}{{01T2}{source}{}{sourcepointer-2}{}}
\newlabel{lemma-locally-finite-type-characterize@cref}{{[originallemma][0][]01T2}{[1][0][]source}{}{}{}}
\makeatother
\crefname{originallemma}{original lemma}{original lemmas}
\Crefname{originallemma}{Original lemma}{Original lemmas}
\makeatletter
\newlabel{lemma-dominant-finite-number-irreducible-components}{{01RM}{source}{}{sourcepointer-3}{}}
\newlabel{lemma-dominant-finite-number-irreducible-components@cref}{{[originallemma][0][]01RM}{[1][0][]source}{}{}{}}
\makeatother
\crefname{originallemma}{original lemma}{original lemmas}
\Crefname{originallemma}{Original lemma}{Original lemmas}
\makeatletter
\newlabel{lemma-dimension-fibre-at-a-point}{{02FX}{source}{}{sourcepointer-4}{}}
\newlabel{lemma-dimension-fibre-at-a-point@cref}{{[originallemma][0][]02FX}{[1][0][]source}{}{}{}}
\makeatother
\crefname{originallemma}{original lemma}{original lemmas}
\Crefname{originallemma}{Original lemma}{Original lemmas}
\makeatletter
\newlabel{algebra-lemma-integral-over-field}{{00GS}{source}{}{sourcepointer-5}{}}
\newlabel{algebra-lemma-integral-over-field@cref}{{[originallemma][0][]00GS}{[1][0][]source}{}{}{}}
\makeatother
\crefname{originallemma}{original lemma}{original lemmas}
\Crefname{originallemma}{Original lemma}{Original lemmas}
\makeatletter
\newlabel{lemma-finite-over-dense-open}{{0BAI}{source}{}{sourcepointer-6}{}}
\newlabel{lemma-finite-over-dense-open@cref}{{[originallemma][0][]0BAI}{[1][0][]source}{}{}{}}
\makeatother
\title{Characterisations of a dominant morphism of finite degree}
\author{The Stacks Project Authors}
\date{stacks-project:morphisms.tex:local-snapshot-20261012; GNU FDL 1.2 or later}
\begin{document}
\maketitle
\setcounter{subsection}{1}
\noindent\textit{Editorial note.} All delivered mathematical text is the author's own, copied from the pinned Stacks Project chapter file morphisms.tex; no statement or proof step was rewritten, reordered or removed.\par\medskip
\noindent\textit{Editorial note.} The chapter file pinned-original-source.tex carries only \textbackslash{}input\{preamble\} on its first line, so no preamble text lives inside the chapter. The retained file sources/stacks-morphisms/retained-preamble.tex is a mechanical, line-for-line copy of the project's single shared preamble as staged in the snapshot. Nothing was added, removed or redefined, and the excerpt serves only as preamble, never as a span coordinate.\par\medskip
\noindent\textit{Editorial note.} Both counted spans start at the \textbackslash{}begin line of their environment, so the statement and the author's proof are delivered whole; surrounding prose of the section is deliberately outside the counted spans and is carried as standard prerequisites instead.\par\medskip
\noindent\textit{Editorial note.} The citation fields (source\_id, source\_version, authors, source\_url, license.name, source\_title) are identical to those of every other card this lane has produced from this chapter, so the batch selector treats them as one source.\par\medskip
\noindent\textit{Allowed standard prerequisites (editorial).}\begin{itemize}
\item Finiteness of a finite type extension of domains detected by the field extension, up to localisation. Algebra, Lemma \textbackslash{}ref\{algebra-lemma-finite-extension-finitely-generated-fractions\} and the surrounding discussion, used at proof lines 13418-13445. Translates between the field and ring statements.
\item Generic point conditions for finite type maps of integral schemes. Lemma \textbackslash{}ref\{lemma-generically-finite\} and Lemma \textbackslash{}ref\{lemma-quasi-finiteness-over-generic-point\}, used at proof lines 13446-13459. Give the pointwise formulations.
\end{itemize}\medskip
\ifdefined\hypertarget\hypertarget{sourcepointer-1}{}\fi
\noindent\textit{Source reference.} Original lemma 02NX: focused-seeds/AGP/sources/stacks-snapshot-20261011/tags.txt, source lines 3496--3496. The lemma's own canonical tag, quoted so that readers can resolve the label inside the delivered text; it points at the counted statement itself and carries no proof content.\par\medskip
\ifdefined\hypertarget\hypertarget{sourcepointer-2}{}\fi
\noindent\textit{Source reference.} Original lemma 01T2: focused-seeds/AGP/sources/stacks-snapshot-20261011/tags.txt, source lines 2359--2359. Navigation pointer to a statement the counted proof refers to but does not contain; nothing load-bearing is delegated to it.\par\medskip
\ifdefined\hypertarget\hypertarget{sourcepointer-3}{}\fi
\noindent\textit{Source reference.} Original lemma 01RM: focused-seeds/AGP/sources/stacks-snapshot-20261011/tags.txt, source lines 2307--2307. Navigation pointer to a statement the counted proof refers to but does not contain; nothing load-bearing is delegated to it.\par\medskip
\ifdefined\hypertarget\hypertarget{sourcepointer-4}{}\fi
\noindent\textit{Source reference.} Original lemma 02FX: focused-seeds/AGP/sources/stacks-snapshot-20261011/tags.txt, source lines 3208--3208. Navigation pointer to a statement the counted proof refers to but does not contain; nothing load-bearing is delegated to it.\par\medskip
\ifdefined\hypertarget\hypertarget{sourcepointer-5}{}\fi
\noindent\textit{Source reference.} Original lemma 00GS: focused-seeds/AGP/sources/stacks-snapshot-20261011/tags.txt, source lines 618--618. Navigation pointer to a statement the counted proof refers to but does not contain; nothing load-bearing is delegated to it.\par\medskip
\ifdefined\hypertarget\hypertarget{sourcepointer-6}{}\fi
\noindent\textit{Source reference.} Original lemma 0BAI: focused-seeds/AGP/sources/stacks-snapshot-20261011/tags.txt, source lines 14211--14211. Navigation pointer to a statement the counted proof refers to but does not contain; nothing load-bearing is delegated to it.\par\medskip
\begin{lemma}
\label{lemma-finite-degree}
Let $X$, $Y$ be integral schemes.
Let $f : X \to Y$ be locally of finite type.
Assume $f$ is dominant.
The following are equivalent:
\begin{enumerate}
\item the extension $R(Y) \subset R(X)$ has
transcendence degree $0$,
\item the extension $R(Y) \subset R(X)$ is finite,
\item there exist nonempty affine opens $U \subset X$
and $V \subset Y$ such that $f(U) \subset V$
and $f|_U : U \to V$ is finite, and
\item the generic point of $X$ is the only point of $X$ mapping to
the generic point of $Y$.
\end{enumerate}
If $f$ is separated or if $f$ is quasi-compact, then these are
also equivalent to
\begin{enumerate}
\item[(5)] there exists a nonempty affine open $V \subset Y$ such
that $f^{-1}(V) \to V$ is finite.
\end{enumerate}
\end{lemma}

\begin{proof}
Choose any affine opens $\Spec(A) = U \subset X$
and $\Spec(R) = V \subset Y$ such that $f(U) \subset V$.
Then $R$ and $A$ are domains by definition. The ring map
$R \to A$ is of finite type
(Lemma \ref{lemma-locally-finite-type-characterize}).
By Lemma \ref{lemma-dominant-finite-number-irreducible-components}
the generic point of $X$ maps to the generic point of $Y$
hence $R \to A$ is injective. Let $K = R(Y)$ be the fraction field
of $R$ and $L = R(X)$ the fraction field of $A$. Then $L/K$
is a finitely generated field extension. Hence we see that
(1) is equivalent to (2).

\medskip\noindent
Suppose (2) holds. Let $x_1, \ldots, x_n \in A$ be generators
of $A$ over $R$. By assumption there exist nonzero polynomials
$P_i(X) \in R[X]$ such that $P_i(x_i) = 0$. Let $f_i \in R$ be the
leading coefficient of $P_i$. Then we conclude that
$R_{f_1 \ldots f_n} \to A_{f_1 \ldots f_n}$ is finite, i.e., (3) holds.
Note that (3) implies (2). So now we see that (1), (2) and (3) are all
equivalent.

\medskip\noindent
Let $\eta$ be the generic point of $X$, and let $\eta' \in Y$ be the
generic point of $Y$. Assume (4). Then
$\dim_\eta(X_{\eta'}) = 0$ and we see that $R(X) = \kappa(\eta)$ has
transcendence degree $0$ over $R(Y) = \kappa(\eta')$ by
Lemma \ref{lemma-dimension-fibre-at-a-point}.
In other words (1) holds. Assume the equivalent conditions (1), (2) and
(3). Suppose that $x \in X$ is a point mapping to $\eta'$.
As $x$ is a specialization of $\eta$,
this gives inclusions $R(Y) \subset \mathcal{O}_{X, x} \subset R(X)$,
which implies $\mathcal{O}_{X, x}$ is a field, see
Algebra, Lemma \ref{algebra-lemma-integral-over-field}.
Hence $x = \eta$. Thus we see that (1) -- (4)
are all equivalent.

\medskip\noindent
It is clear that (5) implies (3) with no additional assumptions on
$f$. What remains is to prove that if $f$ is either separated or
quasi-compact, then the equivalent conditions (1) -- (4) imply (5).
This follows from Lemma \ref{lemma-finite-over-dense-open}.
\end{proof}

\end{document}
